\section{Related work}\label{sec:related-work}

Free completions of partial algebras and partial-to-total encodings are
classical; the construction recalled in \cref{sec:free-completion} is used only
as background \cite{ClimentVidalCosmeLlopez2023,Diaconescu2009}.

The finite-subset theorem is naturally compared with residuality and
discrimination.  In the standard terminology, discrimination asks for a single
homomorphism that is injective on a prescribed finite subset.  What is special
in \cref{thm:finite-subset-discrimination} is the restriction on the codomain:
the whole original carrier, possibly infinite, must remain pointwise fixed and
only finitely many elements may lie outside it.  The usual finite-product proof
need not stay inside this class, as shown by
\cref{prop:finite-complement-product-obstruction}.

For finite bases, the construction lies close to familiar finite-recognition
and residual-finiteness arguments for term algebras and tree automata
\cite{Brainerd1969,Kozen1977,ComonEtAl2008}.  Selected finite subterm sets and a
sink state are standard ingredients.  The present theorem does not claim those
ingredients as new; it isolates their use under a pointwise-fixed possibly
infinite base.

Finite-complement ideas also occur in semigroup theory, notably through Rees
index and related notions \cite{RuskucThomas1998,CainMaltcev2013}.  The analogy
is not literal here because, when $D$ is genuinely partial, the retained copy
of $A$ need not be a subalgebra of the total free completion.  Accordingly we
use the descriptive phrase \emph{finite-complement completion} rather than Rees
index terminology.

The completion in \cref{sec:completion} is also distinct from ordinary
profinite completion.  Profinite constructions use finite quotients or finite
target algebras \cite{DaveyGouveiaHaviarPriestley2011,ChenAdamekMiliusUrbat2016}.
Here the target may itself be infinite because the base is retained pointwise;
only the complement is finite.  The Cauchy-completion machinery and the use of
finite periodic dynamics are standard.  The result of interest is the
properness consequence under this different finiteness constraint.

The arithmetic examples belong to a large literature on partial and totalized
division \cite{Carlstrom2004,BergstraHirshfeldTucker2009,BergstraTucker2022}.
They are examples only and are not used to claim a new treatment of division by
zero.

The narrow claim submitted for specialist evaluation is therefore the direct
finite-subset discrimination theorem for pointwise-base-preserving
finite-complement completions, together with the completion consequences proved
from that same restricted class.  No priority claim is made here beyond that
precise formulation.
